\pdfmapfile{+cm.map}
\pdfmapfile{+cmextra.map}
\pdfmapfile{+symbols.map}
\documentclass[10pt]{article}
\usepackage[margin=0.85in]{geometry}
\usepackage{amsmath,amssymb,amsthm,enumitem,hyperref,fancyhdr}
\hypersetup{colorlinks=true,urlcolor=blue,linkcolor=blue}
\pagestyle{fancy}\fancyhf{}\lhead{K215: direct seeds and shared normalization}\rhead{22 October 2026}\cfoot{\thepage}
\setlength{\headheight}{14pt}\setlength{\parindent}{0pt}\setlength{\parskip}{6pt}\emergencystretch=4em
\setlist[itemize]{label=$\bullet$,leftmargin=*,nosep}\setlist[enumerate]{leftmargin=*,nosep}
\newcommand{\code}[1]{\texttt{#1}}
\newtheorem{proposition}{Proposition}
\begin{document}
\begin{center}{\Large\bfseries Use a smaller seed theory, not another quotient}\\[6pt]
Further Rocq and Lean proof simplification for ProveIt\\22 October 2026\end{center}
\textbf{Recommendation.} In Lean, instantiate completeness with only the original named seed bundles, derive full seed semantics on that same model, and bypass the first indiscernibility quotient. In Rocq, that direct interpreted-seed route already exists: simplify its duplicated bundle normalization instead of removing a quotient it does not have. This report makes the asymmetry explicit and gives a smaller proposed freshness interface.

\textbf{Review boundary.} Main is pinned to
\code{57541fe0cfe3af62b4cb58747d8d1d7c16b338d5}, 22 October, 01:00 UTC.
The review inspected source, commit history and upstream documentation. No proof assistant, incoming verifier or generator was run. Proposed changes below are not compiled patches; no machine or repository source was changed.

\section*{1. Recent changes: adoption is not the same as merge history}
The maintained phase-twelve Lean checkpoint already separates \code{ZFModels} and \code{WFModelCore}, retains compatibility wrappers and removes the unwanted enumerator import from the L2 library. Its documented endpoint closure is \textbf{179 project modules}, down from 192; the endpoint-plus-decoder union is 184, with 15 freshly compiled and 169 source-bound reused providers. The separate 67-module compatibility scope does not make this a fresh 199-module build. These are inspected upstream records, not our replay.

The latest merges preserve those checked sources. The later enumerator-corollary relocation is retained as a proposal outside the selected endpoint after the broader cut; it is not another additive endpoint saving. Shared implication-basis alternatives are also distinguished from the retained phase-eleven core in the maintained checkpoint. Proposing these already catalogued changes as new would obscure the actual next obligation.

The maintained Rocq lane still has the older 42-module source boundary with the offset certificate. Historical normalized-source receipts do not certify different maintained files. This review does not propose a new source-certificate migration.

\textbf{What the new review adds.} Source 186 already proposes a direct full seed for Lean. The current Rocq \code{Complete222.v} already supplies a closely corresponding interpreted-seed proof spine. The useful next step is therefore a carefully scoped Lean port, strengthened by (i) a narrow seed-only theory and (ii) one structural bound for freshness. Separately, Rocq can share its soundness/completeness bundle normalization at the pure named-logic level.

\section*{2. The important asymmetry between the two assistants}
Lean's current path is:
\[
\text{model of DerivableNF}\ \Longrightarrow\ \text{equality-free seed}\
\Longrightarrow\ \text{indiscernibility quotient}\ \Longrightarrow\ \text{FullSeed}.
\]
In \code{KEndpoint.c222\_of\_zf\_refutation}, the proof calls
\code{fullSeed\_of\_derivable}, obtains \code{FullSeed (qmem mem)}, and passes the witness
\code{KIndisc222.mk mem (e 0)} to the lower model construction.
The original completeness continuation, however, already supplies satisfaction of the
\emph{original named formulas}, including their literal equality, not merely satisfaction of their equality translations.
\newpage
Rocq's \code{Complete222.v} already exploits the interpreted predicate
\[
Q(a) := \operatorname{thm222}(\tau(\operatorname{th}(a))).
\]
Its sequence is \code{tau\_fr}, \code{th\_fr}, \code{seed\_matrix},
\code{seed\_code}, \code{seed\_code\_fresh}, \code{Q\_seed},
\code{named\_seed\_to\_seeded}, then \code{TB\_kp0} and the ordinary K0 lower chain.
It does \emph{not} use Lean's first semantic indiscernibility quotient. Rocq's
\code{L224E} is the syntactic equality interpretation needed by \code{Q\_of\_lprov}; it must not be deleted as though it were that quotient.

\section*{3. A narrow direct-seed completeness instance for Lean}
For each arbitrary named formula $\psi$, put
\[
p_\psi=\neg\neg\psi,\qquad c_\psi=\operatorname{enc}(\operatorname{th}(p_\psi)),
\qquad B_\psi=\operatorname{bundleNF}(0,1,c_\psi,p_\psi).
\]
Use the theory
\[
\operatorname{SeedAx}(f)\quad:\!\!\Longleftrightarrow\quad \exists\psi,\ f=B_\psi
\]
in \code{KEind222.c222\_complete}, rather than the theory of all formulas with derivable interpreted codes. This is an interface reduction, not a change to the numerical calculus: each required interpreted seed axiom is still a literal enabled native bundle.

\textbf{Four obligations.}
\begin{enumerate}
\item \textbf{Syntactic commutation.} The bundle's fixed atoms are membership atoms, so
$\operatorname{th}(\operatorname{bundleNF}(a,b,c,p))$
is the same named bundle with matrix $\operatorname{th}(p)$.
Applying \code{enc\_bundleNF} at $c=\operatorname{enc}(\operatorname{th}(p))$ gives the native bundle code. No semantic invariance under equality translation is assumed.
\item \textbf{Enabled gate and freshness.} Double negation makes $c_\psi>1$, so the public-role inequalities hold. Prove that the container label $c_\psi$ does not occur freely in $p_\psi$; the next section gives a potentially smaller sufficient lemma.
\item \textbf{Original named semantics.} The model continuation satisfies $B_\psi$ itself because $\operatorname{SeedAx}(B_\psi)$ holds. Prove a named-bundle normalization theorem for an arbitrary fresh matrix, yielding the seed property for the relation defined by $p_\psi$.
\item \textbf{Remove padding semantically.} Classical double-negation equivalence and \code{seedProp\_congr} give the seed property for $\psi$. As $\psi$ is arbitrary, this is \code{FullSeed mem}, on the original domain.
\end{enumerate}

\textbf{What disappears from this route.} There is no need to restrict the matrix to \code{EqFree}, choose representatives of quotient parameters, prove the satisfaction translation \code{sat\_mk}, or transport the seed witnesses and subset clauses through \code{qmem}. The decoded/named round trip used by the old seed extraction is also unnecessary for this direct extraction. These declarations can remain in legacy modules; their removal from the selected dependency path requires a caller audit.

\textbf{Crucial logical distinction.} The argument uses
$\operatorname{Sat}(f)$ supplied by completeness together with derivability of
$\operatorname{enc}(\operatorname{th}(f))$. It never uses the generally unjustified statement
$\operatorname{Sat}(\operatorname{th}(f))\leftrightarrow\operatorname{Sat}(f)$ on an arbitrary original model. For example, two distinct elements with empty membership are membership-indiscernible, so the translated equality is true while their literal equality is false. This example is not claimed to satisfy the seed theory.
\newpage
\section*{4. One structural bound can replace two freshness inductions}
The Rocq proof separately proves preservation of free names by the equality interpretation and a bound on free names by the encoding. A Lean port need not copy both intermediate interfaces.

The smallest useful single invariant is
\[
\operatorname{Free}(i,f)\ \Longrightarrow\ i<\operatorname{enc}(\operatorname{th}(f)).
\]
Applying it to $i=c_\psi$ and $f=p_\psi$ would give $c_\psi<c_\psi$ if the private label were free. Hence $\neg\operatorname{Free}(c_\psi,p_\psi)$. The existing \code{Tarski.sat\_ext} then removes the container update from the matrix; a second semantic-invariance induction is unnecessary.

\textbf{Proposed proof.} Induct on the original named formula. Membership uses the coordinate bounds for the pairing code. Falsity has no free names. Implication and universal quantification use \code{le\_P\_left}/\code{le\_P\_right} and the encoding factors; a free name under a binder remains free in its body. At equality, expand only the fixed \code{Eind} template: its auxiliary binder $x+y+1$ differs from both original names, which remain in membership atoms. Keep pairing arithmetic opaque except for the established component bounds.

A stronger alternative is $\operatorname{top}(f)\leq\operatorname{enc}(\operatorname{th}(f))$, followed by the existing \code{not\_occurs\_of\_top\_le}. It also handles bound labels, which the seed argument does not require. Prefer the smaller free-name contract unless the stronger bound has independent clients.

These are mathematical proof contracts, not compiled lemmas. Compare the direct proof with Rocq's two small inductions before claiming a net size saving. In particular, \code{sat\_encode\_decode} alone cannot establish freshness of the original formula from its translation.

\textbf{Minimal integration.}
\begin{itemize}
\item Put \code{FullSeed} in a neutral seed-interface module, preserving its qualified name via a wrapper if needed. Currently \code{KLowerSeed} imports \code{KIndisc222} for this interface.
\item Keep \code{formSeed\_of\_fullSeed} unchanged. Its actual argument becomes the direct full seed, and the endpoint witness becomes $e(0)$ instead of the quotient representative $\operatorname{mk}(e(0))$.
\item Keep irreflexivity and the explicit valuation/nonempty-domain witness. Keep the later tagged-domain construction, bisimulation quotient, and well-founded-part construction. This removes only the \emph{first} quotient.
\end{itemize}
The 151-line \code{KIndisc222} is a candidate dependency removal, not a forecast of 151 net lines saved: the new named-bundle and support proofs have costs, and the old API may remain for compatibility. Nor does removal of one quotient imply removal of \code{Quot.sound} or classical choice from the complete endpoint.

\section*{5. Rocq: normalize the bundle once, at the lowest level}
\code{Sound222.v} proves \code{E\_dn}, \code{E\_req}, \code{E\_clos}, \code{E\_seed} for decoded validity. \code{Complete222.v} proves a parallel \code{L\_dn}, \code{L\_req}, \code{L\_clos}, \code{L\_seed}, with a double-negated matrix. Their common content is pure named-formula equivalence.

Prove one family parametrically in a global named equivalence:
\[
\operatorname{leqv}(p,q)\Longrightarrow
\operatorname{leqv}(\operatorname{SeedN}(a,b,c,p),
                    \operatorname{SeedY}(a,b,c,q)).
\]
Analogous request and closure lemmas are internal steps of the same normalization proof.
\newpage
\textbf{Proof and instances.} Expand the request and closure templates. Lift the supplied equivalence through implication, negation and universal quantification using the existing \code{kleqv} congruence lemmas; discharge the fixed double-negation/conjunction normalizations once. The seed theorem is their composition.
\begin{itemize}
\item Soundness: instantiate $p=q$ and reflexivity, then transport the two pure implications through \code{V\_lprov} to obtain the existing \code{E}-equivalence interface.
\item Completeness: instantiate $p=\neg\neg q$ and the existing pure double-negation equivalence. The current \code{L\_seed} is a corollary.
\end{itemize}
No freshness, encoding, irreflexivity, or $a\ne b$ is needed by this \emph{pure normalization} lemma. Those hypotheses are still required where applicable in the outer decoding and seed-to-model adapters. This is theorem-level equivalence, not equivalence derived under arbitrary temporary assumptions; universal congruence must not smuggle in an invalid generalization rule.

\textbf{Dependency layout.} Move the shared named bundle syntax and normalization to a lower module imported by both soundness and completeness. Do not make \code{Sound222} import \code{Complete222}, which already imports soundness. Retain the old lemma names as wrappers until downstream clients and the maintained generator/import lists are checked.

\textbf{Why this is a useful next step.} It removes duplicated normalization arguments and provides a precise cross-assistant template for the Lean named-seed port. It does not remove the separate \code{named\_seed\_to\_seeded} binding/freshness argument or Rocq's syntactic equality interpretation. Net source size and checking time remain unmeasured.

\section*{6. Acceptance order}
\begin{enumerate}
\item Check the pure Rocq normalization module and its two old-interface wrappers. Compile both \code{Sound222} and \code{Complete222}; print their endpoint assumptions and perform recursive checking on the resulting exact sources.
\item In Lean first prove the named-bundle interpretation and one suitable freshness lemma, then derive direct \code{FullSeed} from the narrow \code{SeedAx} model continuation. Test equality-containing matrices, not just equality-free examples.
\item Change the lower seed import and endpoint construction only after the direct theorem compiles. Preserve the explicit witness, ordinary consistency notion, lower tagged quotient, and the unchanged binary machine/program.
\item Rebuild the source-transitively affected endpoint closure and compatibility clients. Report fresh and reused providers separately and audit generator/source correspondence. Do not transfer phase-twelve receipts to modified sources.
\item Measure handwritten proof size, selected dependency size, proof objects and checking time separately. Choose between the combined and split freshness proofs on actual results.
\end{enumerate}

\section*{Sources and attribution}
\small
All current-source references are pinned to \href{https://github.com/VladimirReshetnikov/ProveIt/commit/57541fe0cfe3af62b4cb58747d8d1d7c16b338d5}{57541fe0cfe3}. The direct Lean full-seed idea is credited to Source 186 in \href{https://github.com/VladimirReshetnikov/ProveIt/blob/57541fe0cfe3af62b4cb58747d8d1d7c16b338d5/SetTheory/Cardinals/docs/reports/set-theory/zfc-closure-kernels/intake-213.md}{intake 213}; the narrower seed-only interface combines earlier seed-interface work with that route. The combined free-name bound and shared Rocq normalization are proposals developed in this review.
\begin{itemize}
\item Lean directory \path{Computability/BusyBeaver/Lean/BusyBeaver/ZFC/}: \code{KEind222.lean}, \code{KSeedModel222.lean}, \code{KIndisc222.lean}, \code{KLowerSeed.lean}, \code{KEndpoint.lean}, \code{KLowerQHyp.lean}, \code{TarskiLeibniz.lean}.
\item Rocq directory \path{Computability/BusyBeaver/Research/ZFC/k225-formal/rocq/rocq-extension/coq/215/}: \code{Complete222.v}, \code{Sound222.v}.
\item \href{https://github.com/VladimirReshetnikov/ProveIt/blob/57541fe0cfe3af62b4cb58747d8d1d7c16b338d5/Computability/BusyBeaver/Research/ZFC/k225-formal/verification/k215/simplification/phase12/README.md}{Phase-twelve checkpoint}; \href{https://github.com/VladimirReshetnikov/ProveIt/commit/57541fe0cfe3af62b4cb58747d8d1d7c16b338d5}{latest maintained/proposal merge distinction}.
\end{itemize}
\end{document}
